
In this work, we have shown that the RIT problem is in {\coNP} under GRH. We also showed that $2$-RIT is in {\coRP} assuming GRH, and is in {\coNP} unconditionally. Our algorithms  work by reducing the polynomials modulo a "small" prime~$p$, and taking the computation to the finite field $\FF_p$. We opt for this approach as the coefficients of polynomials represented by algebraic circuits could be doubly exponential in the size of the circuit, which presents problems for an approach via
numerical approximation. We note though the latter approach has been shown to work for deciding  Bounded-RIT  when the exponents~$d_i$ of the radical input~$\mysqrt{d_i}{a_i}$ are unary. In particular, \cite{blomer98} gives a randomised polynomial time algorithm that places  Bounded-RIT  with unary exponents in {\coRP}. 
If we reuse the same exact approach for  Bounded-RIT  with radical inputs with binary exponents, the algorithm no longer runs in polynomial time (the increase in complexity appears in Steps $1$ and $2$ of the algorithm in~\cite{blomer98}, when trying to compute the minimal $d_{ij}$ such that $\mysqrt{d_i}{m_j}^{d_{ij}}\in\ZZ$). Recall our reduction of RIT to its variant where the minimal polynomials of the input radicals $\mysqrt{d_i}{a_i}$ are $x^{d_i}-a_i$ and   all radicands~$a_i$ are pairwise coprime numbers, presented in~\Cref{appendix:reduction}. Applying this reduction first would 
avoid the above-mentioned increase in complexity for  Bounded-RIT  with binary exponents.
It thus places the general version of  Bounded-RIT  in 
{\coRP}. 


\paragraph{\bf Working in extensions of finite fields.}
Our approach to RIT involves
finding a prime $p$ for which the required radicals 
exist in $\FF_p$.  Here it is tempting to consider working instead in a finite extension of $\FF_p$.  For example, it is well-known that every irreducible polynomial over $\FF_p$ of degree $d$ splits completely over $\FF_{p^d}$.
For solving RIT, a major problem with this approach is that if $d$ is given in binary then representing an element of the field $\FF_{p^d}$ requires space
exponential in the bit-length of $d$.  Specialising to 2-RIT, we have that the 
required square roots all exist in $\FF_{p^2}$, for any $p$.  But working in $\FF_{p^2}$ is only sound if the latter is a quotient of the number field $K$ generated by the square roots, that is, if $p$ has inertial degree~2 over $K$.  Moreover, the asymptotic density of such primes is the same as for those that split over $K$.  
